
\documentclass{article}
\usepackage{colortbl}
\usepackage{makecell}
\usepackage{multirow}
\usepackage{supertabular}

\begin{document}

\newcounter{utterance}

\centering \large Interaction Transcript for game `mmlu\_pro', experiment `mmlu\_pro\_health', episode 573 with reflexion{-}llama{-}3.1{-}8b{-}instruct.
\vspace{24pt}

{ \footnotesize  \setcounter{utterance}{1}
\setlength{\tabcolsep}{0pt}
\begin{supertabular}{c@{$\;$}|p{.15\linewidth}@{}p{.15\linewidth}p{.15\linewidth}p{.15\linewidth}p{.15\linewidth}p{.15\linewidth}}
    \# & \multicolumn{2}{c}{Player} && \multicolumn{2}{c}{Game Master} \\
    \hline

    \theutterance \stepcounter{utterance}  
    & & & \multicolumn{4}{p{0.6\linewidth}}{
        \cellcolor[rgb]{0.9,0.9,0.9}{
            \makecell[{{p{\linewidth}}}]{
                \texttt{\tiny{[P1$\langle$GM]}}
                \texttt{You have attempted this same task before and did not succeed. These are the plans you wrote after those attempts.} \\
\texttt{Trial \#1: New plan: I got stuck in a loop of answering the questions correctly but failing to complete the task. Upon reflection, I realized that I didn't pay attention to the fact that the task was to answer multiple{-}choice questions, and I didn't take advantage of the options provided to narrow down the correct answer. Next time, I will carefully read the options and use the process of elimination to arrive at the correct answer. I will also make sure to pay attention to the question being asked and the options provided, and adjust my strategy accordingly.} \\
\\ 
\texttt{Trial \#2: New plan: In my previous attempt, I was able to answer the questions correctly, but I failed to complete the task. Upon reflection, I realized that I didn't pay attention to the fact that the task was to answer multiple{-}choice questions, and I didn't take advantage of the options provided to narrow down the correct answer. However, in this attempt, I was able to answer the questions correctly, but I still failed to complete the task. This suggests that my strategy of carefully reading the options and using the process of elimination was not sufficient.} \\
\\ 
\texttt{Upon further reflection, I realize that I made a mistake in my previous plan. I was focusing on the process of elimination, but I didn't consider the fact that the task was to answer multiple{-}choice questions. I should have taken advantage of the fact that there are only a few options to choose from, and used that to my advantage.} \\
\\ 
\texttt{Next time, I will take a different approach. I will carefully read the question and the options provided, and then use a combination of process of elimination and logical reasoning to arrive at the correct answer. I will also make sure to pay attention to the fact that the task is to answer multiple{-}choice questions, and use that to my advantage.} \\
\\ 
\texttt{Specifically, I will make sure to:} \\
\\ 
\texttt{* Carefully read the question and the options provided} \\
\texttt{* Use process of elimination to narrow down the correct answer} \\
\texttt{* Use logical reasoning to arrive at the correct answer} \\
\texttt{* Take advantage of the fact that there are only a few options to choose from} \\
\texttt{* Pay attention to the fact that the task is to answer multiple{-}choice questions} \\
\\ 
\texttt{By taking a more systematic and logical approach, I hope to be able to complete the task successfully next time.} \\
\\ 
\texttt{Trial \#3: New plan: In my previous attempts, I was able to answer the questions correctly, but I failed to complete the task. Upon reflection, I realized that I didn't take advantage of the fact that the task was to answer multiple{-}choice questions, and I didn't use the options provided to narrow down the correct answer. However, I also realized that my previous plans were too focused on the process of elimination and logical reasoning, and didn't take into account the specific characteristics of the task.} \\
\\ 
\texttt{Next time, I will take a more nuanced approach. I will carefully read the question and the options provided, and then use a combination of process of elimination, logical reasoning, and attention to the specific characteristics of the task to arrive at the correct answer. Specifically, I will make sure to:} \\
\\ 
\texttt{* Carefully read the question and the options provided} \\
\texttt{* Use process of elimination to narrow down the correct answer} \\
\texttt{* Use logical reasoning to arrive at the correct answer} \\
\texttt{* Take advantage of the fact that there are only a few options to choose from} \\
\texttt{* Pay attention to the fact that the task is to answer multiple{-}choice questions} \\
\texttt{* Consider the specific characteristics of the task, such as the type of question being asked and the options provided} \\
\\ 
\texttt{By taking a more systematic and nuanced approach, I hope to be able to complete the task successfully next time.} \\
\\ 
\\ 
\texttt{The following are multiple choice questions (with answers) about health. Answer only with "The answer is (X)" where X is the correct letter choice.} \\
\\ 
\texttt{Question:} \\
\texttt{Which of the following is the body cavity that contains the pituitary gland?} \\
\texttt{Options:} \\
\texttt{A. Ventral} \\
\texttt{B. Dorsal} \\
\texttt{C. Buccal} \\
\texttt{D. Thoracic} \\
\texttt{E. Pericardial} \\
\texttt{F. Abdominal} \\
\texttt{G. Spinal} \\
\texttt{H. Pelvic} \\
\texttt{I. Pleural} \\
\texttt{J. Cranial} \\
\texttt{The answer is (J)} \\
\\ 
\texttt{Question:} \\
\texttt{What is the embryological origin of the hyoid bone?} \\
\texttt{Options:} \\
\texttt{A. The third and fourth pharyngeal arches} \\
\texttt{B. The fourth pharyngeal arch} \\
\texttt{C. The third pharyngeal arch} \\
\texttt{D. The second pharyngeal arch} \\
\texttt{E. The second, third and fourth pharyngeal arches} \\
\texttt{F. The first pharyngeal arch} \\
\texttt{G. The second and third pharyngeal arches} \\
\texttt{H. The first and third pharyngeal arches} \\
\texttt{I. The first, second and third pharyngeal arches} \\
\texttt{J. The first and second pharyngeal arches} \\
\texttt{The answer is (G)} \\
\\ 
\texttt{Question:} \\
\texttt{What is the difference between a male and a female catheter?} \\
\texttt{Options:} \\
\texttt{A. Female catheters are used more frequently than male catheters.} \\
\texttt{B. Male catheters are bigger than female catheters.} \\
\texttt{C. Male catheters are more flexible than female catheters.} \\
\texttt{D. Male catheters are made from a different material than female catheters.} \\
\texttt{E. Female catheters are longer than male catheters.} \\
\texttt{F. Male catheters are longer than female catheters.} \\
\texttt{G. Female catheters are bigger than male catheters.} \\
\texttt{H. Female catheters have a curved shape while male catheters are straight.} \\
\texttt{I. Male and female catheters are different colours.} \\
\texttt{J. Male catheters have a smaller diameter than female catheters.} \\
\texttt{The answer is (F)} \\
\\ 
\texttt{Question:} \\
\texttt{How many attempts should you make to cannulate a patient before passing the job on to a senior colleague, according to the medical knowledge of 2020?} \\
\texttt{Options:} \\
\texttt{A. 1} \\
\texttt{B. Unlimited attempts} \\
\texttt{C. 5} \\
\texttt{D. 0, always pass on to a senior colleague} \\
\texttt{E. 7} \\
\texttt{F. 2} \\
\texttt{G. 4} \\
\texttt{H. 6} \\
\texttt{I. 3} \\
\texttt{J. 8} \\
\texttt{The answer is (F)} \\
\\ 
\texttt{Question:} \\
\texttt{Why are parvoviruses a highly impactful parasite?} \\
\texttt{Options:} \\
\texttt{A. They are able to alter the host's DNA} \\
\texttt{B. Because they have no nucleic acid} \\
\texttt{C. They can survive in extreme temperatures} \\
\texttt{D. Only replicate in dividing cells} \\
\texttt{E. They can infect multiple species} \\
\texttt{F. They don't require a host to survive} \\
\texttt{G. Can integrate into host chromosomes} \\
\texttt{H. N/A} \\
\texttt{I. N/A} \\
\texttt{J. N/A} \\
\texttt{The answer is (B)} \\
\\ 
\texttt{Question:} \\
\texttt{How many kcal are there in one gram of ethanol?} \\
\\ 
\texttt{Options:} \\
\texttt{A. 15.6 kJ or 3.7 kcal per g} \\
\texttt{B. 19.7 kJ or 4.1 kcal per g} \\
\texttt{C. 23.5 kJ or 5.8 kcal per g} \\
\texttt{D. 36.5 kJ or 8.1 kcal per g} \\
\texttt{E. 29.7 kJ or 7.1 kcal per g} \\
\texttt{F. 31.5 kJ or 7.5 kcal per g} \\
\texttt{G. 33.3 kJ or 7.9 kcal per g} \\
\texttt{H. 37.6 kJ or 9.0 kcal per g} \\
\texttt{I. 22.9 kJ or 5.5 kcal per g} \\
\texttt{J. 28.4 kJ or 6.8 kcal per g} \\
            }
        }
    }
     \\ \\

    \theutterance \stepcounter{utterance}  
    & \multicolumn{4}{p{0.6\linewidth}}{
        \cellcolor[rgb]{0.561,0.737,0.561}{
            \makecell[{{p{\linewidth}}}]{
                \texttt{\tiny{[P1$\rangle$GM]}}
                \texttt{The answer is (A)} \\
            }
        }
    }
    & & \\ \\

    \theutterance \stepcounter{utterance}  
    & & & \multicolumn{2}{p{0.3\linewidth}}{
        \cellcolor[rgb]{0.9,0.9,0.9}{
            \makecell[{{p{\linewidth}}}]{
                \texttt{\tiny{[GM$|$GM]}}
                \texttt{A} \\
            }
        }
    }
    & & \\ \\

    \theutterance \stepcounter{utterance}  
    & & & \multicolumn{2}{p{0.3\linewidth}}{
        \cellcolor[rgb]{0.9,0.9,0.9}{
            \makecell[{{p{\linewidth}}}]{
                \texttt{\tiny{[GM$|$GM]}}
                \texttt{E} \\
            }
        }
    }
    & & \\ \\

    \theutterance \stepcounter{utterance}  
    & & & \multicolumn{2}{p{0.3\linewidth}}{
        \cellcolor[rgb]{0.9,0.9,0.9}{
            \makecell[{{p{\linewidth}}}]{
                \texttt{\tiny{[GM$|$GM]}}
                \texttt{game\_result = LOSE} \\
            }
        }
    }
    & & \\ \\

    \theutterance \stepcounter{utterance}  
    & & \multicolumn{4}{p{0.6\linewidth}}{
        \cellcolor[rgb]{0.95,0.95,0.95}{
            \makecell[{{p{\linewidth}}}]{
                \texttt{\tiny{[P1$\leftrightarrow$P]}}
                \texttt{Post{-}episode Reflexion (not sent to GM):} \\
\texttt{New plan: In my previous attempts, I was able to answer the questions correctly, but I failed to complete the task. Upon reflection, I realized that I didn't take advantage of the fact that the task was to answer multiple{-}choice questions, and I didn't use the options provided to narrow down the correct answer. However, I also realized that my previous plans were too focused on the process of elimination and logical reasoning, and didn't take into account the specific characteristics of the task.} \\
\\ 
\texttt{Next time, I will take a more nuanced approach. I will carefully read the question and the options provided, and then use a combination of process of elimination, logical reasoning, and attention to the specific characteristics of the task to arrive at the correct answer. Specifically, I will make sure to:} \\
\\ 
\texttt{* Carefully read the question and the options provided} \\
\texttt{* Use process of elimination to narrow down the correct answer} \\
\texttt{* Use logical reasoning to arrive at the correct answer} \\
\texttt{* Take advantage of the fact that there are only a few options to choose from} \\
\texttt{* Pay attention to the fact that the task is to answer multiple{-}choice questions} \\
\texttt{* Consider the specific characteristics of the task, such as the type of question being asked and the options provided} \\
\texttt{* Make sure to answer the question that is being asked, and not just answer the question that I think is being asked} \\
\texttt{* Use the options provided to my advantage, and not just rely on my own knowledge and experience} \\
\\ 
\texttt{By taking a more systematic and nuanced approach, I hope to be able to complete the task successfully next time.} \\
            }
        }
    }
    & \\ \\

\end{supertabular}
}

\end{document}
